\begin{pastexam}[label={px:ncku-113-cce-12}]{成大電通所 113 學年度 簡答題第 3 題}
（簡答）

(6 pts) Floyd-Warhsall algorithm uses the dynamic programming to find all pair shortest paths.
Let $A^k(i,j)$ denote the value of the shortest path from node $v_i$ to node $v_i$ according to
a vertex whose index is $k$. Please show the recursive function of $A^k(i,j)$ and give the time
complexity of the algorithm.
\end{pastexam}

\begin{pastsoln}
原題有兩個筆誤：``Warhsall'' 是 Warshall；``from node $v_i$ to node $v_i$'' 的終點應該是
$v_j$。 ``according to a vertex whose index is $k$'' 照 Floyd--Warshall 的標準定義解讀成
「中間點的編號都不超過 $k$」。

\textbf{定義}。 頂點是 $v_1, \ldots, v_n$。 $A^k(i, j)$ 是從 $v_i$ 到 $v_j$、中間經過的頂點
（不含兩端）都在 $\{v_1, \ldots, v_k\}$ 裡的最短路徑長度。 $k = 0$ 表示不能經過任何中間點，
只能直接走邊。

\textbf{遞迴式}（資結書第~18.9~節、第~25.6~節）：
\[
A^0(i, j) = \mathit{cost}(i, j) =
\begin{cases}
0, & i = j,\\
w(v_i, v_j), & i \ne j \text{ 且有邊 } (v_i, v_j),\\
\infty, & \text{其他},
\end{cases}
\]
\[
A^k(i, j) = \min\bigl( A^{k-1}(i, j),\ A^{k-1}(i, k) + A^{k-1}(k, j) \bigr),
\qquad 1 \le k \le n .
\]
答案是 $A^n(i, j)$。

\textbf{理由}。 看中間點都在 $\{v_1, \ldots, v_k\}$ 裡的一條最短路徑 $p$（沒有負環時可以取
簡單路徑）：
\begin{itemize}
\item $p$ 不經過 $v_k$：它的中間點都在 $\{v_1, \ldots, v_{k-1}\}$，長度是 $A^{k-1}(i, j)$。
\item $p$ 經過 $v_k$（簡單路徑只經過一次）：從 $v_k$ 切成 $v_i \leadsto v_k$ 和
$v_k \leadsto v_j$ 兩段，兩段的中間點都在 $\{v_1, \ldots, v_{k-1}\}$，各自也必須最短（否則
換成更短的那段，整條會更短），長度是 $A^{k-1}(i, k) + A^{k-1}(k, j)$。
\end{itemize}
兩種情形取小的，就是遞迴式。

\textbf{時間複雜度}：$\Theta(n^3)$。 $k$ 從 $1$ 到 $n$，每個 $k$ 要算 $n^2$ 個 $(i, j)$，每個
$O(1)$。 空間可以只用一個 $n \times n$ 陣列原地更新，$\Theta(n^2)$（第 $k$ 輪的
$A(i, k)$、$A(k, j)$ 在這一輪不會變，因為沒有負環時 $A^k(i, k) = A^{k-1}(i, k)$）。
\end{pastsoln}
